% !TeX root = ../papers/nc-p3.tex
\section{Walsh三次元族の全点に対するAS正則性の必要十分条件}
\label{sec:walsh-normal-quadrics}
\begin{theorem}\label{thm:walsh-normal-quadrics}
標数零の代数閉体でWalsh基底を持つ隣接階数6の四窓を取る。
標準14641型quadratic Koszul左右AS正則であるための必要十分条件は、両切断の
非零色が三つずつあり、欠色を3にした符号付き平方行列 $N$ の階数が2であることである。
従って反対階数16を含む全ての特殊値で、この判定に高次の経路計算は要らない。
正則点では四つの正規二次式が欠色の語を張り、その二乗は共通の四中心四次式を与える。
各始点でそれらの多項式環上自由階数256となる。
\end{theorem}
\begin{proof}[階数1の排除]
欠色が二つある側は定理\ref{thm:walsh-symmetric-windows}の成長障害で除かれる。
$N$ は行和・列和が零なので、残る可能性は階数1,2だけである。
階数1で三行三列が非零なら全九成分が非零で、係数の平方は行倍率と列倍率の積になる。
平方根を取ると $m_{be}=u_b v_e$ の九成分の符号変更となり、
$u_0^2=u_1^2+u_2^2$, $v_0^2=v_1^2+v_2^2$ である。
重み行列の符号判定により $m=(u_bv_e)$ とその $(0,1)$ 成分だけを反転したものの二表示で覆われる。
両円錐曲線を $(1+t^2,1-t^2,2t)$ で有理表示し、二独立変数の有理関数体上で
四次に元の六関係を各位置で挿入する。余核の文字次元は両表示で $(12,8,8,8)$、
合計36である。\path{computations/verify_walsh_normal_quadrics.py} が有理関数体の
正確な階数を再現する。両既約族の全特殊化でも四次元数は36以上であり、標準値35に反する。
\end{proof}

\begin{proof}[欠色四語を一文字の左右で移す写像]
以後 $\operatorname{rank}N=2$ とする。$y_{p,i}=x_{p,i}x_{p+1,i+3}$ と書く。
任意の文字 $a$ に対し
\[
 y_{p,i}x_{p+2,a}=\sum_j T_{p,a}[j,i]\,x_{p,a}y_{p+1,j+3}
 \tag{*}\label{eq:walsh-square-transport}
\]
となる可逆行列 $T_{p,a}$ が一意に定まることを示す。
$h=a+3$ の三文字を $(h+e,h+b+e,h+b)$ と書く。元の三次経路の双対は
$4\times4$ の成分 $X_{be}$ のうち各非欠色行・列が対応する関係直交平面に入るものである。
その次元は5である。実際各側の関係挿入が6次元、交叉が
定理\ref{thm:walsh-dual-frobenius}で1次元、挿入前は16次元である。
両側の欠色成分が零である部分は一次元であり、$W$ の最後の文字を固定した縮約に
命題\ref{prop:walsh-dual-square-cut}の対角文字倍を施した $C$ が張る。

第3行全体を零にしたとき、各非欠色列は $X_{be}=c_e C_{be}$ に制限される。
各非欠色行を第3列まで延長する条件は $\sum_e N_{be}c_e=0$ と同値である。
実際その関係直交平面は $s_e$ を対角とする形式の全等方平面で、
第3列を忘れた射影は単射である。像は $C_{b,*}$ との重み付き直交条件で与えられる二平面である。
$C_{be}^2=m_{be}^2$（巡回すると転置の可能性もある）なので所定の式となる。
階数2では $c_e$ は全て同じで、延長は一意である。この議論は個別の成分では割らない。
従って第3行への制限は全射で核が $kC$ である。第3列への制限も同じ核を持つ全射である。
これを双対化すれば式\eqref{eq:walsh-square-transport}の両端は同じ四空間の二基底である。

全等方性から、各 $X$ で非欠色の行ごとに $\sum_e s_e X_{be}^2=0$、
非欠色の列ごとに $\sum_b s_b X_{be}^2=0$ である。
$\sum_{be}s_bs_e X_{be}^2$ を二通りに取れば二つの境界の二次形式は一致する。
境界の順序はいずれも文字に $a+3$ を加えたものなので、共通符号を除き
$D=\operatorname{diag}(1,-1,-1,1)$ を得る。$D^{-1}=D$ により $T_{p,a}$ もその等長写像である。
$T_p=T_{p,0}$ は両側に共通な語 $x_{p,0}x_{p+1,3}x_{p+2,0}$ を固定するので、
第0基底ベクトルを固定する。行列式は $-1$ である。これは階数2の既約開集合上の
局所正則な関数で二値しか取らないから、一点で決まる。
前掲コードは $m$ の三行を $(0,1,i),(i,0,i),(-1,1,0)$ とした点で
$T_0=\operatorname{diag}(1,1,-1,1)$、全始点で行列式 $-1$ を得る。
\end{proof}

\begin{proof}[四つの正規二次式]
四語の係数ベクトルを分裂四元数環 $M_2(k)$ の基底 $e_0=1,e_1,e_2,e_3=e_1e_2$ で表す。
$e_a e_i=(-1)^{a_1 i_0}e_{a+i}$ とし、余因子共役を $z\mapsto\bar z$ とすれば
行列式ノルムは上の $D$ である。$T_p$ は1を固定する向き反転等長写像なので
$T_p(z)=q_p\bar z q_p^{-1}$ と書ける。実際、余因子共役を合成したものは
跡零三空間の特殊直交群の元であり、$\mathrm{PGL}_2\simeq\mathrm{SO}_3$ の随伴作用で尽くされる。
全頂点のラベルを $a$ だけずらすWalsh作用を式\eqref{eq:walsh-square-transport}に施すと
$T_{p,a}=L_{e_a}T_p L_{e_a}^{-1}$ となる。左右の基底を $i,j$ とした符号比は
$(-1)^{a_1(i+j)_0}$ である。ここで $L$ は左乗法を表す。

次の位置との整合を確認する。三語 $(i,j,k)$ を逆順にして
$(-1)^{b_0 e_1}$ を掛けると、位置 $p$ の三次から位置 $p+1$ の三次への線型同型となる。
確かに辺色 $c$ の関係空間は文字を $c$ ずらして第1文字 $t$ に $(-1)^{c_1t_0}$ を
掛ける操作で保たれる。二位置移動後の各関係は両文字への $J(i)=(-1)^{i_0i_1}$ 倍で得られる。
$h=i+j+k$ を固定すると $b_0e_1$ は、それぞれ $e_1 k_0$ と $e_1h_0$ だけ、
$b_1j_0+i_0i_1+j_0j_1$ と $b_0h_1$ だけ異なる。
これは各辺を逆にする際の必要な符号と、辺色だけの無害な倍率である。
境界の二基底を逆順にすれば
\[
 T_{p+1}=\theta_{\rm hi}\,T_p^{-1}\theta_{\rm lo},
 \qquad q_{p+1}=e_1q_p^{-1}e_2
\]
と選べる。ここで $\theta_{\rm hi}$ と $\theta_{\rm lo}$ は該当ビットの文字倍、
すなわち $e_1,e_2$ の各共役である。二回適用すると $q_{p+2}=e_3q_pe_3$、
四回で元に戻る。

係数ベクトル $e_c q_p$ の定める二次式を $u_{p,c}$ とする。式\eqref{eq:walsh-square-transport}
で次の位置の実際の係数順に戻す移動 $j\mapsto j+3$ は $z\mapsto e_1 z e_2$ である。
四元数の積則から
\[
 u_{p,c}x_{p+2,a}=\kappa_c\det(q_p)\,\phi_c(a)x_{p,a}u_{p+1,c},
 \quad
 \phi_c(a)=(-1)^{[a\ne0]+a_0c_1+a_1c_0},\quad
 \kappa_c=(-1)^{[c\ne0]+c_1}.
\]
実際 $q_p\overline{e_c q_p}=\det(q_p)\bar e_c$ を使い、
$e_a\bar e_c e_a/\det(e_a)=\phi_c(a)\bar e_c$ とすればよい。
$\det(q_{p+1})=\det(q_p)^{-1}$ より四位置で定数倍率の積は1である。
$u$ を位置ごとに規格化して係数を $\phi_c(a)$ のみにできる。
各 $q_p$ が可逆なので、四式は各位置の欠色空間の基底である。

折り畳んだ環で $u_c=\sum_p u_{p,c}$ と取る。矢を二位置ずらし
$\phi_c$ を掛ける対合で交換するため $u_c^2$ は中心である。
また $u_c$ と $u_d$ は非零スカラーで位置ごとに交換する。
実際 $\phi_c(i)\phi_c(i+3)$ は $s_i$ に一定符号を掛けたもので、
係数ベクトルへの $s_i$ 倍は $e_3$ の共役である。
従って $e_d q_p$ の直線は $e_d q_{p+2}$ の直線に移る。
\end{proof}

\begin{proof}[全次数の自由性とKoszul性]
四つの $u$ を全て割った商は定理\ref{thm:walsh-dual-frobenius}の
$(1+t)^4$ の窓である。各始点でその基底語を持ち上げる。
次数帰納と上の交換式により、順序付けた $u_c$ の単項式と持上げ語で全経路を張る。
各 $u_c^2$ を中心変数に移せば、中心四変数上の生成系の母関数の上界は
$(1+t)^4(1+t^2)^4=(1+t+t^2+t^3)^4$ となる。
各開始位置で $u$ の個数に応じて語の開始位置が移るが、各位置の商基底の次数数は同じである。
中心変数を含めた上界は $(1-t)^{-4}$ である。
一方この既約族の幾何学的一般点は正則なので、各次数の経路次元は全特殊化で
その標準値以上である。全次数で上下界が一致し、形式的中心変数上の自由加群からの
全射は同型である。特に中心四変数は独立である。
有限Frobenius双対と標準Hilbert級数に加え、中心多項式環上自由であるから、
定理\ref{thm:klein-rational-period-four}の左右線型複体の証明が適用できる。
両側で中心の最初の二変数が正則なので単純から環への $\operatorname{Hom},\operatorname{Ext}^1$
が消え、表示の完全性とその双対で両端を、Euler等式で中央を得る。
従って全次数でKoszul左右AS正則となる。
\end{proof}
